% Adapted from docs/wiki/Installation.md; edit this file, not the wiki page, for the manual.
\chapter{Installation}\label{ch:installation}

Exacta requires Python 3.10 or newer, PyTorch 2.4 or newer and NumPy 1.22 or newer.
Install PyTorch for your hardware using its \href{https://pytorch.org/get-started/locally/}{installation guide},
then install Exacta. Pin the release you evaluated, so that results can be traced to it:

\begin{lstlisting}[language=bash]
python -m pip install exacta==X.Y.Z   # X.Y.Z: the release you use
python -c "import exacta; print(exacta.__version__)"
\end{lstlisting}

The wheels are for Linux x86-64 with glibc 2.28 or newer. Python 3.10 and 3.11
have dedicated wheels; the Python 3.12 stable-ABI wheel supports Python 3.12 and newer.

Other platforms install from source, which builds the CPU extension and requires a C++17
compiler and Python development headers; OpenMP is optional. On Debian/Ubuntu, install
\texttt{build-essential} and development headers matching your Python interpreter.

\section{CPU and GPU}\label{installation:cpu-and-gpu}

CPU kernels are precompiled in wheels and compile during source installation. CUDA loss and
metric kernels use Triton at runtime; their first invocation can include compilation. The CPU extension does not link against a
specific PyTorch binary. Local validation covers Linux x86-64 CPU and NVIDIA CUDA execution;
other hardware/platform combinations should be verified before use.

\section{Distributed CPU execution}\label{installation:distributed-cpu-execution}

For Gloo distributed runs, use PyTorch 2.8 or newer. PyTorch 2.4.1 can deadlock during
process-group destruction after collecting metric scores: a Gloo worker needs the Python
GIL while the outer destructor holds it. PyTorch's \href{https://github.com/pytorch/pytorch/releases/tag/v2.8.0}{2.8 release notes}
include the destructor fix (\#154976). Ordinary non-distributed CPU evaluation still supports
PyTorch 2.4.

\section{Optional references}\label{installation:optional-references}

SciPy is needed for Hungarian lesion matching. MONAI, MedPy, VTK, SimpleITK and other reference
libraries are used by optional validation tests; they are not needed for ordinary Dice or
HD95 evaluation. Chapter~\ref{ch:tests} lists the packages and how to install them.

\section{Development checkout}\label{installation:development-checkout}

Use \texttt{main} to evaluate merged changes beyond the published release. Create a feature branch
from it when contributing. This manual describes the version printed on its title page; a newer
checkout may expose additional APIs or changed behavior. To keep your own Python environment
untouched, create a virtual environment first (Section~\ref{tests:referee}).

\begin{lstlisting}[language=bash]
git clone --branch main https://github.com/rodolfo-viana/exacta.git
cd exacta
python -m pip install -e ".[test]"
\end{lstlisting}

Record the commit SHA and environment alongside experiment results. Reinstall after changing
C++ sources so the compiled extension matches the checkout.
